% Generated by experiments/scripts/export_latex.py -- do not edit by hand.
% Source: experiments/data/2-processed-tables-and-evaluations/revised-usd-12m/robustness/robustness.json
% Source: experiments/data/2-processed-tables-and-evaluations/revised-usd-12m/benchmark/benchmark.json
\begin{table}[htbp]
\centering
\caption[평가 집합에 따른 지표 가족 평균의 민감도.]{평가 집합이 바뀌면 지표 가족 평균이 얼마나 움직이는지 보여 줘요. 매칭 코호트에서 SHA256으로 늘 같은 결과가 나오게 뽑은 부분 표본마다(시드 \texttt{contract-benchmark-v1}, 코호트 전체는 계약 $n=6{,}483$개) 엔도스랭크(ER)와 AWP의 지표 가족 평균 켄달의 $\tau$를 구했어요. 단계마다 그 단계의 계약에 닿는 화살표로 점수를 매겼어요. 열 이름에서 “허락”은 허락, “송금”은 송금, “시빌”은 시빌 안정성을 뜻해요. 행마다 평가한 계약 집합이 달라요. 그래서 이 $\tau$ 값들은 표~\ref{tab:alignment-hybrid}의 값과 비교할 수 없어요. 이 값들은 평가하는 집단에 따라 통계량이 얼마나 움직이는지 보여 줘요. 단계마다 구간은 계산하지 않았어요. 실행 시간은 표~\ref{tab:benchmark-scaling}의 측정값이에요.}
\label{tab:robustness-sample-size}
\fitwidth{%
\begin{tabular}{rccccccrrr}
\toprule
$n$ & ER 허락 & AWP 허락 & ER 송금 & AWP 송금 & ER 시빌 & AWP 시빌 & ER (초) & AWP (초) & AWP/ER 시간 \\
\midrule
1{,}000 & 0.665 & 0.228 & 0.225 & 0.752 & 0.049 & 0.441 & 0.498 & 1.307 & 2.6$\times$ \\
2{,}000 & 0.611 & 0.223 & 0.200 & 0.731 & 0.036 & 0.421 & 0.318 & 2.383 & 7.5$\times$ \\
5{,}000 & 0.577 & 0.229 & 0.201 & 0.735 & 0.033 & 0.420 & 1.672 & 6.620 & 4.0$\times$ \\
6{,}483 & 0.560 & 0.226 & 0.196 & 0.736 & 0.026 & 0.419 & 2.657 & 8.530 & 3.2$\times$ \\
\bottomrule
\end{tabular}
}
\end{table}
